% ============================================================
%  4. Proposed Method
% ============================================================
\section{Proposed Method}
\label{sec:method}

\subsection{Overview}

\LightRet{} is trained via a three-stage progressive distillation pipeline:
\[
  \underbrace{\textsc{BERT}}_{\text{Teacher}_1}
  \;\xrightarrow{\;\text{Stage 1}\;}
  \underbrace{\RetBERT{}}_{\text{Teacher}_2}
  \;\xrightarrow{\;\text{Stage 2}\;}
  \underbrace{\LightRet{}}_{\text{Backbone}}
  \;\xrightarrow{\;\text{Stage 3}\;}
  \underbrace{\LightRet{}\text{-NER}}_{\text{Final model}}
\]
Each stage uses a frozen teacher to supervise a smaller student via a
task-specific distillation loss, bridging the representational gap between the
full BERT encoder and the final compact model progressively.

% ─────────────────────────────────────────────────
\subsection{Stage 1: Sentence-Level Distillation
            (\textsc{BERT}\,$\to$\,\RetBERT{})}
\label{subsec:stage1}

\paragraph{Teacher.}
BERT-base-uncased~\citep{devlin2019bert} (110M params, all frozen).
Its sentence representation is the mean of the final-layer hidden states:
\begin{equation}
  \bz^{B} \;=\; \frac{1}{|\bw|}\sum_{i=1}^{|\bw|} \bh^{B}_{i}
  \;\in\; \RR^{768}.
  \label{eq:bert_sent}
\end{equation}
\noindent where $\bz^{B} \in \RR^{768}$ is the BERT sentence vector, $|\bw|$ is the number of input tokens, and $\bh^{B}_{i}$ is the $i$-th final-layer hidden state of BERT-base.

\paragraph{Student: \RetBERT{}.}
\RetBERT{} replaces BERT's WordPiece embedding lookup with a frozen
\RetVec{} embedder followed by a linear projection:
\begin{equation}
  \bh^{(0)}_{i}
  \;=\;
  \bW_{p}\,\RetVec(w_i) \;+\; \text{PE}(i),
  \quad \bW_{p} \in \RR^{768 \times 256},
  \label{eq:retbert_input}
\end{equation}
\noindent where $\bh^{(0)}_{i}$ is the initial hidden state for word $i$, $\bW_{p} \in \RR^{768 \times 256}$ is a trainable linear projection matrix, $\RetVec(w_i) \in \RR^{256}$ is the frozen character-level embedding of word $w_i$, and $\text{PE}(i)$ is the fixed sinusoidal positional encoding
of~\citet{vaswani2017attention}.
The student then applies $L_{RB}=12$ standard pre-LayerNorm Transformer
encoder layers~\citep{xiong2020layer}:
\begin{equation}
  \bh^{(\ell)}_{i}
  \;=\;
  \text{TransformerLayer}_{\,d=768,\;H=12}
  \!\left(\bh^{(\ell-1)}\right)_{i},
  \quad \ell = 1, \ldots, 12.
  \label{eq:retbert_layers}
\end{equation}
\noindent where $\bh^{(\ell)}_{i}$ is the hidden state at Transformer layer $\ell$ for word $i$, $d{=}768$ is the hidden dimension, $H{=}12$ is the number of attention heads, and layers are stacked for $\ell = 1, \ldots, 12$.

Its sentence representation is:
\begin{equation}
  \bz^{RB} \;=\; \frac{1}{n}\sum_{i=1}^{n} \bh^{(12)}_{i}
  \;\in\; \RR^{768}.
  \label{eq:retbert_sent}
\end{equation}
\noindent where $\bz^{RB} \in \RR^{768}$ is the RetBERT sentence vector obtained by mean-pooling the final Transformer layer's hidden states over $n$ words.

\paragraph{Loss.}
We minimise the cosine distance between teacher and student sentence vectors:
\begin{equation}
  \calL_{1}
  \;=\;
  1 - \frac{\bz^{B} \cdot \bz^{RB}}
           {\|\bz^{B}\| \cdot \|\bz^{RB}\|}
  \;\in\; [0,\, 2].
  \label{eq:loss1}
\end{equation}
\noindent where $\cdot$ denotes the dot product, $\|\cdot\|$ the $\ell_2$ norm, $\bz^{B}$ and $\bz^{RB}$ are the teacher and student sentence vectors respectively. A value of $0$ indicates identical directions; $2$ indicates opposite directions.

Optimising $\calL_1$ directly encourages \RetBERT{} to produce
BERT-compatible sentence representations from character-level inputs, without
requiring intermediate layer matching.

% ─────────────────────────────────────────────────
\subsection{Stage 2: Token-Level Compression
            (\RetBERT{}\,$\to$\,\LightRet{})}
\label{subsec:stage2}

\paragraph{Teacher.}
\RetBERT{} (frozen after Stage 1) produces per-token hidden states
$\{\bh^{RB}_{i}\}_{i=1}^{n} \subset \RR^{768}$.

\paragraph{Student: \LightRet{} backbone.}
\LightRet{} is a compact BiGRU--Transformer model operating in a
$d=256$-dimensional space:
\begin{align}
  \be_i &= \RetVec(w_i) \;\in\; \RR^{256},
  \label{eq:lr_embed}\\[4pt]
  \bg_i &= \bigl[
             \overrightarrow{\text{GRU}}_{128}(\be_{1:n})_i\;;\;
             \overleftarrow{\text{GRU}}_{128}(\be_{1:n})_i
           \bigr]
           \;\in\; \RR^{256},
  \label{eq:lr_bigru}\\[4pt]
  \bh^{(\ell)}_i &= \text{TransformerLayer}_{\,d=256,\;H=4}
                    \!\left(\bh^{(\ell-1)}\right)_i,
                    \quad \ell = 1,\ldots,4,
  \label{eq:lr_transformer}
\end{align}
where $\be_i \in \RR^{256}$ is the frozen RetVec embedding of word $w_i$; $\overrightarrow{\text{GRU}}_{128}$ and $\overleftarrow{\text{GRU}}_{128}$ are forward and backward GRU units with 128 hidden units each; $[\,;\,]$ denotes vector concatenation; $\bg_i \in \RR^{256}$ is the BiGRU output; $\bh^{(\ell)}_i$ is the hidden state after the $\ell$-th Transformer layer ($d{=}256$, $H{=}4$ heads); and $\bh^{(0)} \coloneqq \bg$.

For Stage~2 only, a temporary linear projector aligns the student's
output dimension to the teacher's:
\begin{equation}
  \bp_i \;=\; \bW_{\text{proj}}\,\bh^{(4)}_i + \bb_{\text{proj}},
  \quad \bW_{\text{proj}} \in \RR^{768 \times 256}.
  \label{eq:lr_proj}
\end{equation}
\noindent where $\bp_i \in \RR^{768}$ is the projected student token representation aligned to the teacher's 768-d space; $\bW_{\text{proj}} \in \RR^{768 \times 256}$ and $\bb_{\text{proj}} \in \RR^{768}$ are temporary projection weights that are \emph{discarded} before Stage~3.

\paragraph{Loss.}
Token-level cosine distillation with sequence-length masking:
\begin{equation}
  \calL_{2}
  \;=\;
  \frac{1}{n} \sum_{i=1}^{n}
  \left(
    1 - \frac{\bh^{RB}_{i} \cdot \bp_{i}}
             {\|\bh^{RB}_{i}\| \cdot \|\bp_{i}\|}
  \right).
  \label{eq:loss2}
\end{equation}
\noindent where $\bh^{RB}_{i} \in \RR^{768}$ is the frozen RetBERT teacher's hidden state at token position $i$, $\bp_{i} \in \RR^{768}$ is the projected student representation (Eq.~\ref{eq:lr_proj}), and $n$ is the number of tokens in the sentence.

By supervising every token position, Stage~2 teaches \LightRet{} the
\emph{contextual} structure of \RetBERT{}'s hidden space, not just sentence
summaries.

% ─────────────────────────────────────────────────
\subsection{Stage 3: Noisy-Student NER Fine-Tuning}
\label{subsec:stage3}

\subsubsection{Model Setup}

After Stage~2, the projector is removed and a BiLSTM NER head is attached:
\begin{align}
  \bs_i &= \bigl[
             \overrightarrow{\text{LSTM}}_{128}(\bH)_i \;;\;
             \overleftarrow{\text{LSTM}}_{128}(\bH)_i
           \bigr]
           \;\in\; \RR^{256},
  \label{eq:ner_bilstm}\\
  \text{logits}_i &= \bW_{c}\,\bs_i + \bb_{c},
  \quad \bW_{c} \in \RR^{|\calY| \times 256}.
  \label{eq:ner_logits}
\end{align}
where $\bH = \{\bh^{(4)}_i\}_{i=1}^{m}$ is the LightRet backbone output sequence for \emph{noisy} input words; $\bs_i \in \RR^{256}$ is the BiLSTM output at position $i$ formed by concatenating 128-unit forward and backward LSTM states; $\bW_{c} \in \RR^{|\calY| \times 256}$ and $\bb_{c} \in \RR^{|\calY|}$ are the linear classifier weight and bias; and $|\calY|{=}9$ is the number of BIO entity labels.

\subsubsection{Noise Augmentation}

Each clean training sentence $\bw$ is perturbed character-by-character using
the operators of Equations~\eqref{eq:nsub}--\eqref{eq:nspace} with
$p_{\text{sub}}{=}0.10$, $p_{\text{ins}}{=}0.05$,
$p_{\text{del}}{=}0.05$, $p_{\text{space}}{=}0.02$.
A \emph{shift log} records each edit as a triple
$(k, \delta, \tau)$ where $k$ is the character position, $\delta \in \{-1,0,+1\}$
is the length change, and $\tau \in \{\text{sub},\text{ins},\text{del},\text{space}\}$
is the operation type.

\subsubsection{Dynamic BIO Label Projection}
\label{subsec:label_proj}

Because $\calN_{\text{space}}$ can increase the word count ($m > n$), we
cannot naively copy labels from $\bw$ to $\tilde{\bw}$.  We define a
word-level alignment $\calA = \{(C_k, N_k)\}_{k=1}^{K}$ derived from the
shift log, where $C_k \subseteq \{1,\ldots,n\}$ and
$N_k \subseteq \{1,\ldots,m\}$ are the sets of clean and noisy word indices
in the $k$-th group.  Labels are projected as follows:

\begin{equation}
\tilde{y}_j \;=\;
\begin{dcases}
  y_{C_k[0]}                 & \text{if } |C_k|=|N_k|=1 \quad \text{(1:1)},\\
  y_{C_k[0]}                 & \text{if } |N_k|=1 \quad \text{(merge: many clean} \to \text{one noisy)},\\
  y_{C_k[0]}                 & \text{if } j = N_k[0] \quad \text{(split: first fragment)},\\
  \sigma\!\left(y_{C_k[0]}\right) & \text{if } j \in N_k\setminus\{N_k[0]\} \quad \text{(split: continuation fragments)},
\end{dcases}
\label{eq:label_proj}
\end{equation}
\noindent where $\tilde{y}_j$ is the projected BIO label assigned to noisy word $j$; $C_k \subseteq \{1,\ldots,n\}$ and $N_k \subseteq \{1,\ldots,m\}$ are the sets of clean and noisy word indices belonging to alignment group $k$; $y_{C_k[0]}$ is the original label of the first clean word in the group; and the continuation function $\sigma : \calY \to \calY$ maps \texttt{B-X} $\mapsto$ \texttt{I-X} and is the identity on all other labels, ensuring BIO validity is preserved after word splitting.

\subsubsection{Stage 3 Loss}

The teacher (\LightRet{}, frozen, clean input) and student (\LightRet{}
+ NER head, trainable, noisy input) are trained jointly.
Let $\bH^{T} = \{\bh^{T}_{i}\}_{i=1}^{n}$ be the teacher's hidden states
and $\bH^{S} = \{\bh^{S}_{j}\}_{j=1}^{m}$ be the student's.

For alignment-aware distillation, each group $k$ contributes the
cosine distance between the teacher's representation of group $C_k$ and
the student's representation of group $N_k$:
\begin{equation}
  \bh^{T}_{(k)} \;=\; \frac{1}{|C_k|}\sum_{i \in C_k} \bh^{T}_{i},
  \qquad
  \bh^{S}_{(k)} \;=\; \frac{1}{|N_k|}\sum_{j \in N_k} \bh^{S}_{j}.
  \label{eq:align_pool}
\end{equation}
\noindent where $\bh^{T}_{(k)} \in \RR^{256}$ and $\bh^{S}_{(k)} \in \RR^{256}$ are the mean-pooled teacher and student hidden representations for alignment group $k$; $\bh^{T}_{i}$ and $\bh^{S}_{j}$ are the per-token hidden states of the teacher (clean input) and student (noisy input) respectively.

The distillation and classification losses are:
\begin{align}
  \calL_{\text{distill}}
  &= \frac{1}{K} \sum_{k=1}^{K}
     \left(1 - \frac{\bh^{T}_{(k)} \cdot \bh^{S}_{(k)}}
                    {\|\bh^{T}_{(k)}\| \cdot \|\bh^{S}_{(k)}\|}\right),
  \label{eq:loss_distill}\\[6pt]
  \calL_{\text{class}}
  &= -\frac{1}{m} \sum_{j=1}^{m}
     \sum_{c=1}^{|\calY|} \mathbf{1}[\tilde{y}_j = c]
     \log \hat{p}_{j,c},
  \label{eq:loss_class}
\end{align}
\noindent where $K$ is the total number of alignment groups; $\calL_{\text{distill}}$ is the alignment-aware cosine distillation loss; $m$ is the number of noisy words; $\tilde{y}_j$ is the projected label for noisy word $j$ (Eq.~\ref{eq:label_proj}); $\mathbf{1}[\cdot]$ is the indicator function; and $\hat{p}_{j,c} = \text{softmax}(\text{logits}_j)_c$ is the predicted probability for BIO label $c$ at position $j$.

The combined Stage~3 loss is:
\begin{equation}
  \calL_{3} \;=\; \beta\,\calL_{\text{class}}
               \;+\; (1-\beta)\,\calL_{\text{distill}},
  \quad \beta = 0.5.
  \label{eq:loss3}
\end{equation}
\noindent where $\beta \in [0,1]$ is a scalar interpolation coefficient balancing the classification signal ($\calL_{\text{class}}$) and the teacher-alignment signal ($\calL_{\text{distill}}$). We set $\beta{=}0.5$ in all experiments, confirmed via a validation sweep (Section~\ref{sec:analysis}).
